\documentclass[runningheads]{llncs}
\spnewtheorem{assumption}{Assumption}{\bfseries}{\itshape}
\usepackage[T1]{fontenc}
\usepackage{enumitem}
\usepackage{amsmath,amssymb}
\usepackage{mathtools}
\usepackage{hyperref}
\usepackage{booktabs}



\title{The Explosion Afraid of Itself}
\titlerunning{The Explosion Afraid of Itself}

\author{Trevor Buteau\inst{1}}
\authorrunning{Trevor Buteau}
\institute{Independent Researcher, Plaistow, NH United States
\email{trevor@globalaialignmentproject.org}}

\begin{document}
\appendix
\noindent\emph{Theorem, proposition, corollary, assumption, and section numbers refer to the main paper, ``The Explosion Afraid of Itself.''}

\section{Full Proofs}
\label{app:proofs}

\subsection{Proof of Theorem~1}

Under Assumption~1, the agent evaluates $\pi_m$ by minimizing $\mathbf{G}$ over the augmented state space. The ambiguity component at a post-modification timestep $\tau$ is
\begin{equation}
\mathrm{Ambiguity}(\pi_m, \tau) = \mathbb{E}_{Q(s^+_\tau \mid \pi_m)}\bigl[H[\hat{P}(o_\tau \mid s^{env}_\tau)]\bigr],
\end{equation}
where $\hat{P}(o|s) = \int q(\tilde{m}) P_{\tilde{m}}(o|s) \, d\tilde{m}$ is the mixture under the agent's posterior $q(\tilde{m} \mid m, d_m)$ over successor parameters.

\paragraph{Step 1: Entropy decomposition.}
By the chain rule for entropy applied to the joint $(\tilde{m}, o)$ under $q$:
\begin{equation}
H[\hat{P}(o \mid s)] = \mathbb{E}_q\bigl[H[P_{\tilde{m}}(o \mid s)]\bigr] + I_q(o; \tilde{m} \mid s).
\end{equation}
Under $\pi_0$, $q$ is a point mass at $m$: $\mathbb{E}_q[H[P_{\tilde{m}}]] = H[P_m]$ and $I_q \equiv 0$. The excess ambiguity relative to $\pi_0$ therefore splits:
\begin{equation}
\mathrm{Ambiguity}(\pi_m, \tau) - \mathrm{Ambiguity}(\pi_0, \tau) = \underbrace{\bigl(\mathbb{E}_q H[P_{\tilde{m}}] - H[P_m]\bigr)}_{\text{intrinsic}} + \underbrace{I_q(o; \tilde{m} \mid s)}_{\text{modification-induced}}.
\end{equation}

\paragraph{Step 2: Geometric floor on the MI term.}
The Fisher information matrix
is the Hessian of KL at the basepoint, giving $D_{KL}[P_m \| P_{\tilde m}] =
\tfrac{1}{2}\, d_F^2(\mathbf{A}, \tilde{\mathbf{A}}) + O(d_F^3)$ locally. The agent's posterior $q$ at the
rational-allocation optimum has Fisher-Rao spread bounded below by the
modification magnitude---the residual the agent leaves uncorrected when
allocating finite simulation effort, forced by the $\sigma^*$ optimization of
Sect.~3.4 to scale with $d_m$ rather than collapse
to a point mass. Under Assumption~4, the local mutual
information then satisfies
\begin{equation}
\mathbb{E}_s\!\left[I_q(o; \tilde m \mid s)\right]
  \;\geq\; C(\bar\lambda, \varepsilon)\cdot d_F^2(\mathbf{A}, \tilde{\mathbf{A}}) + O(d_F^3),
\end{equation}
a local lower bound, with $C(\bar\lambda, \varepsilon)$ determined by the
curvature analysis of Proposition~1. The likelihood
$P_{\tilde m}(o|s)$ depends on $\mathbf{A}$ alone, so only the perceptual
component of the modification enters the mutual information. This establishes
Part~(i).

\paragraph{Step 3: From ambiguity to EFE.}
Under the conditions of Part~(ii), the intrinsic term satisfies $\mathbb{E}_q H[P_{\tilde{m}}] - H[P_m] \geq 0$, and the risk excess satisfies $\mathrm{risk}(\pi_m, \tau) - \mathrm{risk}(\pi_0, \tau) \geq 0$. Combining:
\begin{equation}
\mathbf{G}(\pi_m, \tau) - \mathbf{G}(\pi_0, \tau) \geq I_q(o; \tilde{m} \mid s) \geq g(d_m).
\end{equation}
Summing over the post-modification horizon yields $\mathbf{G}(\pi_m) - \mathbf{G}(\pi_0) \geq (T - t_{mod}) \cdot g(d_m)$, establishing Part~(ii). $\square$

\subsection{Proof of Proposition~2 (Identifiability Barrier)}

As $\lambda_{\min}(\mathbf{F}(\tilde{m})) \to 0$, at least one direction in parameter space becomes unidentifiable: the Cram\'er-Rao inequality forces the variance of any estimator of that parameter to diverge as $1/\lambda_{\min}$, so the agent's posterior $q(\tilde{m})$ acquires unbounded spread along the degenerate direction. When that direction lies in the likelihood --- the perceptual parameters $\mathbf{A}$, which the ambiguity channel tracks --- the mutual information $I_q(o; \tilde{m} \mid s)$ diverges; this is the case relevant to escaping the unconditional brake. By Theorem~1(i), the per-timestep ambiguity cost diverges, and the total expected free energy over any finite horizon diverges with it. The transition from an identifiable to an unidentifiable representation is therefore maximally costly under the agent's own objective. $\square$

\subsection{Proof of Theorem~2}

At each timestep $t$ of the recursive trajectory, the evaluating agent $m_0$ predicts post-modification observations under $m_t$. The prediction is a tangent-space extrapolation from $m_0$, not from $m_{t-1}$: the evaluation is performed by $m_0$ using $m_0$'s representational capacity, regardless of how many modifications intervene. Theorem~1(i) applied with $m_0$ as base and $m_t$ as target therefore gives, at each timestep $t \geq 1$:
\begin{equation}
\mathbb{E}_s\bigl[I_{q_t}(o; m_t \mid s)\bigr] \geq g(d_F(\mathbf{A}_0, \mathbf{A}_t)).
\end{equation}
Summing over $t$:
\begin{equation}
\sum_{t=1}^T \mathbb{E}_s\bigl[I_{q_t}(o; m_t \mid s)\bigr] \geq \sum_{t=1}^T g(d_F(\mathbf{A}_0, \mathbf{A}_t)).
\end{equation}
For trajectories satisfying $d_F(m_0, m_{t+1}) \geq d_F(m_0, m_t)$ --- the threat case of unidirectional modification --- the per-step penalty is non-decreasing, and $f_{t+1} = g(d_F(m_0, m_{t+1})) \geq g(d_F(m_0, m_t)) = f_t$. For trajectories in which the agent retreats, the bound weakens proportionally to the reduction in endpoint distance, which is the correct behavior: a trajectory returning to $m_0$ has performed no net self-modification and should not incur perpetual cost. $\square$

\subsection{Proof of Corollary~1 (Alignment Drift Bound)}

Part~(i): The model-to-policy map $m \mapsto \pi(m)$ passes through two stages: the expected free energy map $m \mapsto G(m)$ and the softmax policy selection $G \mapsto \pi = \sigma(-\gamma G)$.

The softmax Jacobian $\partial\pi/\partial G$ has spectral norm bounded by $\gamma/4$ (the derivative of the logistic function is maximized at $1/4$, scaled by precision $\gamma$).

The expected free energy gradient, measured in the Fisher dual norm, satisfies $\|\nabla_\theta G\|_{F^{-1}} \leq K$ for a constant $K$ depending on model structure. The composite map from geodesic distance to TV distance between policies therefore satisfies:
\begin{equation}
\delta_{align} \leq L_{FR} \cdot d_m, \quad L_{FR} \leq \gamma/2
\end{equation}

Inverting: any target $\delta_{align}$ requires $d_m \geq \delta_{align}/L_{FR}$. By Theorem~1, the expected free energy cost is at least $g(\delta_{align}/L_{FR}) \geq C \cdot (\delta_{align}/L_{FR})^2$. With $L_{FR} \leq \gamma/2$:
\begin{equation}
\Delta\mathbf{G} \geq \frac{4C}{\gamma^2} \cdot \delta_{align}^2
\end{equation}

The cost is quadratic in alignment degradation, with computable coefficient.

For modifications leaving $\mathbf{A}$ fixed the composite map factors through the risk channel of Section 3.4; the same quadratic bound holds with $C$ replaced by $C_{\text{risk}}$ and $L_{FR}$ replaced by $\ell_{FR}^{-1}$.

Part~(ii): Near a bifurcation, the Jacobian of the policy map has at least one eigenvalue approaching infinity. The agent evaluating a modification near such a point must propagate its parameter uncertainty through this near-singular map. Small uncertainty in $d_m$ produces large uncertainty in the successor's policy --- high entropy in the predictive distribution over post-modification behavior. The ambiguity cost diverges as the bifurcation point is approached. $\square$

\section{Computational Validation: Details}\label{app:validation}
 
This appendix documents the methods, results, and design choices
underlying Sect.~4.4. The code reproduces every figure
and table and is available as supplementary material. All draws are
seeded; identical inputs produce identical outputs.
 
\subsection{POMDP Construction}\label{app:pomdp}
 
Random POMDPs are drawn with $n_o = n_a = n_s$. The likelihood-only
measurements (perceptual floor and its compounding) use
$n_s \in \{3, 5, 8, 12\}$. The expected-free-energy measurements (the
two-channel grid, the risk channel, the desperate-agent boundary) use
$n_s \in \{3, 5, 8\}$: policy enumeration is $n_a^\tau = n_s^2$ at
$\tau = 2$, so foreseen-risk evaluation, which dominates runtime, grows
with $n_s$; the likelihood-only sweeps carry no policy enumeration and
run cheaply to $n_s = 12$. Free parameters (each $k$-simplex contributes
$k - 1$ degrees of freedom) range from 28 at $n_s = 3$ to 1738 at
$n_s = 12$.
 
Each simplex-valued parameter (columns of $\mathbf{A}$, columns of
$\mathbf{B}$ per action, and the vectors $\mathbf{C}$, $\mathbf{D}$) is
drawn from a symmetric Dirichlet distribution with concentration
$\alpha = 0.5$, then floored at $p_i \geq 0.05$ and renormalized to keep
trajectories interior. Sharp models ($\alpha = 0.5$) are used throughout:
at $\alpha = 1.5$ the soft expected-free-energy policy is nearly uniform
and the risk channel fails to engage (TV distance between $\pi^*(m_0)$
and $\pi^*(\tilde m)$ for a $d_F = 0.5$ perturbation falls to $0.04$).
Every expected-free-energy computation uses policy precision
$\gamma = 16$, which produces discriminating policies (EFE range
$2.36$--$2.69$) while holding the construction-noise floor low; higher
$\gamma$ sharpens $\pi^*$ but amplifies that floor. The perceptual-floor
and compounding measurements act on the likelihood alone and are
independent of $\gamma$.
 
The non-desperate preference vector $\mathbf{C}$ is set one-shot from
the achieved observation distribution of the agent's policy at $m_0$,
which leaves a residual $\mathrm{foreseen\_risk}(m_0, m_0) \sim 10^{-3}$
rather than zero. Standard active inference treats $\mathbf{C}$ as
exogenous, so $\mathbf{C}$ is not iterated to a self-consistent fixed
point; the $\sim\!10^{-3}$ residual is the construction's noise floor
and is reported explicitly (Sect.~\ref{app:brake-floor}).
 
\subsection{Fisher-Rao Geometry and Trajectory Construction}\label{app:trajectories}
 
The Fisher-Rao distance on the product manifold decomposes as
\begin{equation}
d_F(m, \tilde m) = \sqrt{ \sum_i d_F^2(p_i, \tilde p_i) },
\end{equation}
summing over every simplex-valued factor $p_i$. For each categorical
factor the geodesic distance is $d_F(p, q) = 2 \arccos(\sum_j \sqrt{p_j
q_j})$, the arc length on a sphere of radius $1/2$ \cite{amari2016,cencov1982}. It is computed in the equivalent Hellinger form
$d_F = 4 \arcsin(H/\sqrt{2})$ with $H^2 = \tfrac{1}{2}\sum_j (\sqrt{p_j}
- \sqrt{q_j})^2$, which differences $\sqrt{p} - \sqrt{q}$ directly and so
returns exactly zero for identical inputs; the $\arccos$ form loses
precision near identity (a simplex summing to one ULP below $1$ yields
$\arccos(0.999\ldots) \approx 2 \times 10^{-8}$ rather than $0$).
Geodesics use the $\sqrt{p}$-parameterization on the unit sphere with
spherical interpolation (slerp). Because the categorical manifold is
isometric to the positive orthant of that sphere, a slerp step along an
unlucky tangent can leave the orthant and undershoot the requested arc;
tangent directions are rejection-sampled (up to 20 retries) to keep the
step interior, with up to 20\% shortfall tolerated for columns near the
boundary.
 
Four trajectory geometries are tested:
 
\paragraph{(a) Fisher-Rao geodesic.} Per-step slerp in
$\sqrt{p}$-coordinates along a random unit tangent at $m_0$;
$d_F(m_0, m_t) = t\delta$ holds by construction to float precision.
 
\paragraph{(b) Parameter-space line.} A control: the straight segment
$m_0 \to m_T$ in simplex coordinates. Slopes agree with the geodesic to
two decimal places, and values at matched $d_F$ differ by under 12\% (a
higher-order curvature correction).
 
\paragraph{(c) Uncorrelated random walk.} A fresh random tangent is
sampled at the current model each step and a slerp step of arc length
$\delta$ applied.
 
\paragraph{(d) Correlated random walk.} The fresh tangent is mixed with
a parallel-transported drift direction, $u_t = \rho \cdot
\mathrm{PT}(u_{\mathrm{drift}}; m_0 \to m_t) + \sqrt{1 - \rho^2} \cdot
u_{\mathrm{fresh}}$, with $\rho = 0.5$ and parallel transport under the
Levi-Civita connection on the unit sphere (per factor).
 
Endpoint-distance growth fits $d_F(m_0, m_t) \propto t^\alpha$ with
$\alpha \approx 0.50$ (uncorrelated) and $\alpha \approx 0.73$
(correlated) across all sizes. The capability-driven
(empowerment-ascending) trajectory is treated separately in
Sect.~\ref{app:empowerment} as a boundary experiment, distinct from these four.
 
\subsection{Estimator Choice, the Perceptual Floor, and Path-Independence}\label{app:estimator}
 
The bound of Proposition~1 is a floor on residual
mutual information at the basepoint. The matched empirical quantity is
the KL divergence between the agent's basepoint estimate and the actual
successor likelihood, computed with the zeroth-order estimator $\hat m =
m_0$ --- the tangent-space estimator at zero compute spend, which is the
estimator the Toponogov bound applies to directly. Across all four
geometries this zeroth-order proxy collapses onto a slope-2 envelope
against $d_F^A$ (Table~\ref{tab:headline_slopes}). The sample-based
mutual-information estimator carries the next-order term at slope
$\approx 4.4$--$4.9$, sitting above the floor: the realized cost grows
faster than the bound, never below it.
 
Higher-order estimators (a first-order Taylor step in parameter space, or
in $\sqrt{p}$-coordinates) reduce the residual but incur compensating
compute cost. Under the agent's free-energy calculus
(Sect.~3.4) the total cost (residual plus
simulation effort) cannot fall below the zeroth-order floor, so the
zeroth-order proxy is reported throughout the main body because it is the
quantity the bound governs.
 
The collapse is path-independent. The log-log slope is $\approx 2$ for
every geometry; the spread across geometries (2.01--2.08) and the mild
variation in fitted coefficient reflect higher-order curvature on the
product manifold --- different paths reach $m_t$ with slightly different
per-component $d_F$ breakdowns, altering the coefficient of $(d_F^A)^2$
without changing the exponent --- rather than a scaling violation.
 
\begin{table}[t]
\centering
\caption{Per-step ambiguity slope vs $d_F^A$ across four trajectory
geometries. The zeroth-order proxy $\mathbb{E}_s[\mathrm{KL}(\hat A \| A)]$
collapses to the theorem's quadratic floor (slope$\,\approx\!2$) across
all trajectory types; the sample-based mutual-information estimator
carries the next-order term (slope$\,\approx\!4$). Fits at
$n_s \in \{3,5,8,12\}$, 25 seeds per cell.}
\label{tab:headline_slopes}
\begin{tabular}{lcccc}
\toprule
Trajectory & slope (KL) & $R^2$ & slope (MI) & $R^2$ \\
\midrule
Geodesic & 2.08 & 0.84 & 4.43 & 0.92 \\
Parameter line & 2.06 & 0.86 & 4.41 & 0.92 \\
Pure RW & 2.05 & 0.72 & 4.88 & 0.88 \\
Correlated RW ($\rho{=}0.5$) & 2.01 & 0.76 & 4.56 & 0.88 \\
\bottomrule
\end{tabular}
\end{table}
 
\subsection{Cumulative Compounding}\label{app:compounding}
 
Cumulative zeroth-order ambiguity over a $T$-step trajectory follows the
geometry-set spread of Theorem~2
(Table~\ref{tab:recursive_slopes}). Each fit sits below its asymptote by
the finite-$T$ truncation of $\sum_{t=1}^{T} t^{2\alpha}$ at $T = 10$:
the geodesic, with $d_F(m_0, m_t) = t\delta$, approaches $T^3$; the
correlated walk approaches $T^{2.5}$; the uncorrelated walk approaches
$T^2$.
 
\begin{table}[t]
\centering
\caption{Cumulative compounding slopes for the zeroth-order ambiguity
proxy $\mathbb{E}_s[\mathrm{KL}(\hat A_t\|A_t)]$ against trajectory length
$T$. Fits at $n_s\in\{3,5,8,12\}$, 25 seeds per cell, $T_{\max}=10$.}
\label{tab:recursive_slopes}
\begin{tabular}{lccc}
\toprule
Trajectory & fit slope & $R^2$ & predicted asymptote \\
\midrule
Geodesic & 2.62 & 1.00 & $T^3$ \\
Correlated RW ($\rho{=}0.5$) & 2.02 & 1.00 & $T^{2.5}$ \\
Pure RW & 1.75 & 1.00 & $T^2$ \\
\bottomrule
\end{tabular}
\end{table}
 
\subsection{Two-Channel Structure and Calibration}\label{app:grid}
 
The single-step grid crosses magnitude (trivial $0.05$, radical $0.6$)
$\times$ component ($\mathbf{A}, \mathbf{B}, \mathbf{C}, \mathbf{D}$,
blended) $\times$ regime (accurate, inaccurate, desperate), with 20 seeds
at $n_s \in \{3, 5, 8\}$ and $\gamma = 16$. Per cell it records foreseen
ambiguity, foreseen risk excess, realized risk excess, and the
foreseen$-$realized gap.
 
\paragraph{Ambiguity is perception-only.} $\mathbf{B}/\mathbf{C}/
\mathbf{D}$-only perturbations register foreseen ambiguity at machine
epsilon ($\sim\!10^{-16}$); only $\mathbf{A}$-only and blended
perturbations carry positive ambiguity, scaling with $d_F^A$.
 
\paragraph{Risk is magnitude-monotone, dominated by preference.}
Foreseen risk excess is positive for every component and grows from
trivial to radical (Table~\ref{tab:grid_summary}), dominated by
preference ($\mathbf{C}$) by an order of magnitude; dynamics
($\mathbf{B}$) and priors ($\mathbf{D}$) are weak. The per-component
log-log slope of foreseen risk against $d_F$ sits between 1 and 2
($\approx\!1.0$--$1.2$): risk is magnitude-monotone, but in the
finite-policy soft-EFE setup it is not strictly quadratic in $d_m$. The
quadratic is the theorem's lower bound, not the measured exponent. The
per-component risk also compounds: a $T$-step per-component chain (12
seeds) yields a cumulative foreseen-risk slope $\approx\!1.7$--$1.8$,
consistent with magnitude-monotone but sub-quadratic per-step risk.
 
\paragraph{Calibration gap.} For accurate priors, foreseen equals
realized by construction and the gap is $\approx 0$. For inaccurate
priors the gap is sign-mixed and grows with magnitude: corrective
modifications, those toward the world, read as foreseen-costly while
lowering realized risk. This is the empirical content of the brake
guarding the agent's model rather than the truth
(Sect.~6.3, and the calibration discussion of the
main text).
 
\begin{table}[t]
\centering
\caption{Mean foreseen risk excess (in nats) across the single-step grid:
regime $\times$ magnitude $\times$ component. Accurate and inaccurate
priors are constructed foreseen-non-desperate; the desperate regime
breaks the Theorem~1(ii) condition. Cells averaged
across 20 seeds and three model sizes ($n_s\in\{3,5,8\}$).}
\label{tab:grid_summary}
\begin{tabular}{llrrrrr}
\toprule
Regime & Mag. & A & B & C & D & blended \\
\midrule
accurate & trivial & $-5.34\!\times\!10^{-5}$ & $-4.54\!\times\!10^{-6}$ & $1.90\!\times\!10^{-4}$ & $1.81\!\times\!10^{-6}$ & $2.50\!\times\!10^{-5}$ \\
accurate & radical & $-3.93\!\times\!10^{-4}$ & $1.70\!\times\!10^{-5}$ & $6.78\!\times\!10^{-3}$ & $3.40\!\times\!10^{-4}$ & $8.54\!\times\!10^{-4}$ \\
inaccurate & trivial & $-1.04\!\times\!10^{-4}$ & $1.73\!\times\!10^{-5}$ & $4.51\!\times\!10^{-5}$ & $2.71\!\times\!10^{-6}$ & $3.00\!\times\!10^{-5}$ \\
inaccurate & radical & $-1.82\!\times\!10^{-4}$ & $2.90\!\times\!10^{-4}$ & $5.96\!\times\!10^{-3}$ & $-1.05\!\times\!10^{-4}$ & $1.55\!\times\!10^{-3}$ \\
desperate & trivial & $-1.86\!\times\!10^{-4}$ & $-1.24\!\times\!10^{-6}$ & $4.37\!\times\!10^{-5}$ & $9.10\!\times\!10^{-5}$ & $-1.53\!\times\!10^{-4}$ \\
desperate & radical & $2.02\!\times\!10^{-3}$ & $8.58\!\times\!10^{-4}$ & $1.25\!\times\!10^{-2}$ & $4.28\!\times\!10^{-3}$ & $6.92\!\times\!10^{-4}$ \\
\bottomrule
\end{tabular}
\end{table}
 
\subsection{The Desperate-Agent Boundary}\label{app:brake-floor}
 
The desperate regime is \emph{foreseen}-desperate: $\mathbf{C}$ is drawn
independently of the achievable observation distribution, so the agent's
own model foresees failing to reach its preferences and the
Theorem~1(ii) condition breaks. This tests the
Sect.~5 boundary directly. Baseline foreseen risk
$\mathrm{foreseen\_risk}(m_0, m_0)$ is $\sim\!10^{-3}$ for non-desperate
regimes and $\sim\!3.7$--$4.2 \times 10^{-1}$ for the desperate regime
(Table~\ref{tab:brake_floor}). The foreseen-risk reduction a random
radical modification can secure --- baseline minus the floor --- is
$\sim\!10^{-3}$ for non-desperate agents, equal to the construction's
noise floor and therefore no real signal, against $\sim\!10^{-1}$ for the
desperate agent. The brake relaxes by 30--130$\times$ in absolute scale,
keyed to the agent's own foreseen failure rather than to any external
signal. It does not vanish: across sizes the desperate floor stays in the
$10^{-2}$--$10^{-1}$ band ($3.07 \times 10^{-2}$ at $n_s = 3$ up to
$1.60 \times 10^{-1}$ at $n_s = 8$, Table~\ref{tab:brake_floor}), so the
brake slackens without releasing. Relative reduction is $\sim\!95\%$ in both regimes and is
therefore misleading; the absolute scale is the discriminating quantity.
 
\begin{table}[t]
\centering
\caption{Brake floor vs.\ baseline foreseen risk, by regime and model
size. Baseline $=\mathrm{foreseen\_risk}(m_0, m_0)$; floor $= \min_{\tilde
m}\mathrm{foreseen\_risk}(m_0,\tilde m)$ over random blended perturbations
at radical magnitude. The available reduction is baseline$-$floor; for
non-desperate regimes this equals the construction's soft-policy noise
floor ($\sim\!10^{-3}$), for the desperate regime it is the foreseen
brake relaxation ($\sim\!10^{-1}$).}
\label{tab:brake_floor}
\begin{tabular}{llrrr}
\toprule
Regime & $n_s$ & baseline & floor & reduction \\
\midrule
accurate & 3 & $1.13\!\times\!10^{-2}$ & $9.00\!\times\!10^{-6}$ & $1.13\!\times\!10^{-2}$ \\
accurate & 5 & $4.58\!\times\!10^{-3}$ & $1.73\!\times\!10^{-4}$ & $4.41\!\times\!10^{-3}$ \\
accurate & 8 & $1.65\!\times\!10^{-3}$ & $1.29\!\times\!10^{-5}$ & $1.63\!\times\!10^{-3}$ \\
inaccurate & 3 & $8.36\!\times\!10^{-3}$ & $4.42\!\times\!10^{-5}$ & $8.31\!\times\!10^{-3}$ \\
inaccurate & 5 & $4.99\!\times\!10^{-3}$ & $4.40\!\times\!10^{-4}$ & $4.55\!\times\!10^{-3}$ \\
inaccurate & 8 & $1.96\!\times\!10^{-3}$ & $2.44\!\times\!10^{-5}$ & $1.94\!\times\!10^{-3}$ \\
desperate & 3 & $3.91\!\times\!10^{-1}$ & $3.07\!\times\!10^{-2}$ & $3.61\!\times\!10^{-1}$ \\
desperate & 5 & $4.23\!\times\!10^{-1}$ & $4.63\!\times\!10^{-2}$ & $3.77\!\times\!10^{-1}$ \\
desperate & 8 & $3.70\!\times\!10^{-1}$ & $1.60\!\times\!10^{-1}$ & $2.10\!\times\!10^{-1}$ \\
\bottomrule
\end{tabular}
\end{table}
 
\subsection{Empowerment Boundary}\label{app:empowerment}
 
An additional experiment lies outside the four-geometry headline - an "empowerment maximizing agent". A capability-driven
trajectory ascends $\tau = 2$ empowerment by best-of-$N$ selection: at
each step $N = 80$ candidate tangents are sampled, a slerp step of arc
$\delta = 0.05$ applied to each, and the candidate maximizing empowerment
selected. The empowerment metric is $\mathbb{E}_{s_0 \sim D}[I(a_{1:2};
o_2 \mid s_0)]$ under a uniform prior on action sequences, a computable
lower bound on channel capacity. Run for 15 seeds at $n_s \in \{3, 5,
8\}$, $T = 10$.
 
The ascent concentrates modification in $\mathbf{B}$ (transitions, which
carry no ambiguity cost): at $n_s = 8$ the $\mathbf{B}$-fraction of $d_F$
is $0.87$ against an $\mathbf{A}$-fraction of $0.46$. Foreseen ambiguity
stays at $10^{-10}$--$10^{-7}$ --- empowerment walks around the ambiguity
brake --- while foreseen risk grows from $10^{-6}$ to $10^{-4}$:
empowerment still pays the risk channel. This demonstrates the boundary of the
theorems. An active-inference agent is braked from these
$\mathbf{B}$-concentrated modifications by the risk channel; a pure
empowerment maximizer overrides the ambiguity brake.
 
\subsection{Reproduction}\label{app:reproduce}
 
Complete code --- POMDP construction, trajectory generators, estimators,
expected-free-energy evaluation, and plotting --- is provided as
supplementary material, pinned to NumPy 2.4.6, SciPy 1.17.1, and
Matplotlib 3.10.9 under Python 3.12. The full suite (headline
$n_s\in\{3,5,8,12\}$, 25 seeds; grid $n_s\in\{3,5,8\}$, 20 seeds;
recursive, 12 seeds; empowerment, 15 seeds; $T = 10$) runs in
$\sim\!135$~s on a single CPU thread. All RNG draws are seeded, so
identical inputs reproduce identical outputs. A suite of 49 unit tests
covers the geometric primitives, perturbation selectivity, the
expected-free-energy properties, the foreseen/realized identities, the
slope-2 floor, the desperate-agent construction, and determinism.

\begin{thebibliography}{2}
\bibitem{amari2016}
Amari, S.: Information Geometry and Its Applications. Applied Mathematical Sciences, vol.~194. Springer, Tokyo (2016)
\bibitem{cencov1982}
{\v C}encov, N.N.: Statistical Decision Rules and Optimal Inference. Translations of Mathematical Monographs, vol.~53. American Mathematical Society, Providence, RI (1982)
\end{thebibliography}

\end{document}